\chapter{DESCRIPCIÓN DEL PROYECTO INTRODUCCIÓN}

    \section{Introducción}\label{intro}

    En el mundo desaparecen miles personas al año, por causas cómo desastres naturales, conflictos armados, aumento del envejecimiento de la población y enfermedades de salud mental. En España, en el año 2024 hubo unas 16.000 denuncias. Actualmente siguen activas, solo del año pasado, unas 700. 700 familias que conviven constantemente con la incertidumbre y la angustia de no saber el paradero de su ser querido. \cite{DatosDesaparecidos2024} Quizás es difícil acabar con las causas de las desapariciones, pero sí que podemos intentar prevenir que una persona llegue a encontrarse en esa situación o que el operativo de búsqueda de esa persona sea lo más efectivo posible, sobretodo si se trata de desapariciones involuntarias. 
    
    Aunque existen distintos tipos de desapariciones, como veremos más adelante, todas las personas tienen derecho a su búsqueda y localización. La búsqueda de personas desaparecidas en la actualidad supone un reto y un esfuerzo notable de carácter social y humanitario\cite{garcia2019personas}. Es por ello que si puedo ayudar a mejorar un poquito el mundo gracias a mis conocimientos de Ingeniería Biomédica quiero aportar mi granito de arena.

    La sociedad en la que vivimos ya se mueve frente a esta problemática. Cada vez existen más herramientas tecnológicas que permiten localizar a una persona. Desde drones con visión artificial y sistemas GPS, hasta aplicaciones de marca España como AlertCops que facilita la seguridad ciudadana y pueden ser de gran ayuda en las búsquedas. Aunque todas estas tecnologías salvan vidas y dan esperanza, todavía hay espacios de la búsqueda que son bastante deficientes, cómo el hecho de que la mayoría de estos sistemas se basan en la idea de que la persona desaparecida lleve consigo algún tipo de dispositivo electrónico localizable. Por ello, es vital llenar estos huecos que quedan, para minimizar la gravedad de estos sucesos y favorecer una mayor seguridad de nuestra sociedad y país.

    La ingeniería biomédica combina conocimientos de salud y tecnología para mejorar la vida y la salud de las personas mediante el uso y desarrollo de soluciones innovadoras. Permite desarrollar desde tratamientos o diagnósticos personalizados, hasta dispositivos de asistencia. Pero no solo es eso, la ingeniería biomédica también se ocupa de la gestión y tratamiento seguro y ético de los datos personales y de salud, garantizando la privaciad y confidencialidad de la información. Por todo ello, la ingeniería biomédica puede ser una estrategia poderosa para proteger y salvar vidas.

    Respecto a la problemática que nos ocupa, la de las personas desaparecidas, puede abordarse desde diferentes puntos de vista de la ingeniería biomédica. Desde la creación de dispositivos para geolocalizar y monitorizar a una persona en tiempo real, usando, por ejemplo, su información fisiológica; hasta el uso de algoritmos para identificar patrones para la localización de las personas desaparecidas. Asimismo, la correcta integración de los datos de las desapariciones es clave para mejorar la coordinación de los operativos de búsqueda.

    En este trabajo me centraré en usar mis conocimientos para definir una app que permita gestionar y coordinar a las diferentes personas o equipos que participan en la búsqueda de una persona desaparecida. Un lugar desde donde gestionar toda la información de una búsqueda y donde se recopilen datos para que se pueda trabajar en el futuro con algoritmos predictivos de localización de la persona desaparecida en base al resultado de otras búsquedas y la información de la propia persona, como su información personal o datos médicos. Este último punto en principio será teórico puesto que no puedo acceder realmente a los datos.





    
    \subsection{Motivación}
    \textcolor{red}{REVISAR, DAR OTRA VUELTA}
    
    Desde hace unos pocos meses soy voluntaria en el GREM (Grupo de Rescate Espeleológico y de Montaña) de Burgos. Entre la multitud de tareas que realizan, trabajan en colaboración con profesionales de las Fuerzas y Cuerpos de seguridad del Estado en la búsqueda de personas desaparecidas. La idea de la aplicación surgió por un vídeo suyo de instagram donde se explicaba que cuanta más gente participase en las búsquedas más probabilidades habría de encontrar a la persona con vida.

    La app pretende reunir a la mayor cantidad de personas posibles a esas búsquedas, además de difundir de forma optimizada las alertas de personas desaparecidas que en muchos casos nos llegan a través de personas desaparecidas. Además, se pretende usar la app para la mejora de la gestión de las búsquedas con distintos elementos, con el objetivo de optimizar recursos. 

    Pero sobretodo, se pretende que la app sea una ayuda para los profesionales, los grupos, las asociaciones, los familiares de las personas desaparecidas y los propios desaparecidos. Se trata de optimizar los recursos para poder salvar vidas.

%%%%%%%%%%%%%%%%%%%%%%%%% CONOCIMIENTOS PREVIOS %%%%%%%%%%%%%%%%%%%%%%%%%
    \section{Conocimientos previos I. Contexto}\label{contexto}

    https://interior.gob.es/opencms/pdf/archivos-y-documentacion/documentacion-y-publicaciones/publicaciones-descargables/seguridad-ciudadana/Atencion_proteccion_orientacion_familias_personas_desaparecidas_folleto_web_126190813.pdf
            
    \textcolor{red}{Contexto: Expón brevemente cómo se abordan los casos de personas desaparecidas en la práctica. Describe el flujo convencional, iniciando con la denuncia, la acción inmediata de las autoridades, la movilización de unidades y la difusión a través de medios y redes sociales.}

Antes de adentrarse de lleno en el la solución que se propone en este documento, se debe conocer un poco su contexto, conociendo de forma resumida el proceso que se lleva a cabo ante una desaparición y las instituciones responsables. 

El abordaje de los casos de personas desaparecidas sigue un flujo bastante complejo, estructurado y coordinado entre distintos organismos. Este proceso comienza con la denuncia de la desaparición de una persona ante fuerzas de seguridad del estado como Policía o Guardia Civil. Estas denuncias son normalmente interpuestas por los familiares, testigos o autoridades. A modo informativo, se ha de añadir que para presentar la denuncia de desaparición no es necesario esperar que pase un plazo específico de tiempo, de hecho cuanto antes se denuncie antes se puede actuar.

Tras interponer la denuncia, los responsables, que normalmente son las fuerzas de seguridad del estado recopilan toda la información que pudiera ser de interés en la búsqueda. Ejemplos de esta infomación son su última ubicación, sus características físicas, circunstancias de desaparición o información médica. Un ejemplo de cuestionario es el que podemos observar en la siguiente imagen \textcolor{red}{INCLUIR IMAGEN/TABLA} y en el Anexo C podemos conocer el flujo de preguntas que normalmente se realiza en estas situaciones.

Una vez recopilada la información se movilizan las unidades de búsqueda especializadas, en estas unidades pueden incluirse grupos ajenos a las fuerzas de seguridad del estado como grupos caninos como el GREM \cite{GREM}. Paralelamente, en algunos casos en los que no entorpezca la búsqueda o investigación del caso, la difusión de estas desapariciones en medios de comunicación o redes sociales puede ser crucial para recibir nuevas pistas. Todo este procese se puede complementar con información de bases de datos nacionales e internacionales.

Es importante reseñar que la combinación de tecnología, rapidez y coordinación es vital para resolver los casos de desapariciones con éxito. Para que estas acciones sean efectivas, se requiere la participación de instituciones sólidas, especializadas y bien estructuradas  que faciliten la labor de los profesionales implicados en las búsquedas. Estas entidades no solo aportan recursos técnicos y humanos, sino que también establecen protocolos de actuación, canalizan la información y garantizan una respuesta organizada ante situaciones de emergencia.


Actualmente en España, existe una institución principal que se encarga de todos estos casos de personas desaparecidas, el Centro Nacional de Desaparecidos (CNDES) \cite{CNDES}. Fue creado en 2018 y actúa como un órgano de gestión de coordinación del sistema de personas desaparecidas utilizado por las Fuerzas y Cuerpos de Seguridad del estado. También desarrollan medidas de cooperación entre instituciones y organizaciones, se encargan de analizar y valorar propuestas de familiares de personas desaparecidas, y es el punto de contacto entre cuerpos policiales y servicios públicos. \cite{garcia2019personas}

Por otro lado, también tenemos instituciones especializadas, como la Fundación ANAR \cite{F_ANAR} fundada en 2017 por el ministerio del interior se encarga de menores desaparecidos.
    
%\textcolor{red}{Mencionar brevemente las normativas legales y protocolos internos que orientan la actuación de las autoridades, poniendo de relieve la importancia de la coordinación y el manejo sensible de datos.}
    
%\textcolor{red}{Exponer las deficiencias o limitaciones de los procesos actuales, lo que permitirá justificar la creación de la app como una herramienta que mejora la coordinación, agilidad y seguridad en el manejo de la información.}
    
    

    
    \subsection{Persona desorientada, perdida y desaparecida}

    \textcolor{red}{Definir Persona desorientada y perdida}
    Persona desorientada es una persona que desconoce la posición que ocupa geográficamente. https://dle.rae.es/desorientar

    Una persona perdida es aquella que tiene una idea de donde se encuentra pero desconoce la forma de llegar a su destino. https://dle.rae.es/perdido?m=form

    La RAE define como persona desaparecida aquella que se haya en paradero desconocido. Se define como "Persona desaparecida la persona ausente de su residencia habitual sin motivo conocido o aparente, cuya existencia es motivo de inquietud o bien que su nueva residencia se ignora, dando lugar a la búsqueda en el interés de su propia seguridad y sobre la base del interés familiar o social”.  (Consejo de Europa, 2009)
    
    % https://cndes-web.ses.mir.es/publico/Desaparecidos/Primeros-Pasos.html
    
    % https://rm.coe.int/16806da843

    % https://sosdesaparecidos.es/como-actuar/#1568743348168-892fd5e6-ad94

    Según los datos proporcionados por el ministerio del interior y el Centro Nacional de Desaparecidos la cifra de desaparecidos en 2024 sube hasta 16.147. \cite{DatosDesaparecidos2024}

    De hecho las estadísticas muestran que... (indicar el porcentaje hombres/ mujeres y por edades, si me puedo centrar en las desapariciones involuntarias mejor.)
    
    Las personas desaparecen en la actualidad por multitud de causas cómo puede ser ...
    
    \subsubsection{Clasificación de casos de desapariciones} 
    Aunque existen distintas formas de clasificar las desapariciones, nosotros las clasificaremos con la clasificación actual propuesta desde el Ministerio del Interior las desapariciones se clasifican en\cite{CNDES,FolletoInfo_personasDes}:
        \begin{itemize}
            \item \textbf{Desapariciones Voluntarias}: en el caso en el que una persona decide ausentarse de su residencia habitual por voluntad propia, sin informar a ningún conocido o familiar.
            \item \textbf{Desapariciones Involuntarias}: en el caso en el por accidente, desorientación o distintas situaciones la persona desaparece sin intención de hacerlo. 
            \item \textbf{Desapariciones Forzosas}: en el caso de que la desaparición se deba por acciones deliberadas, como es el caso de secuestros, violencia o conflictos armados. Es un crimen de la humanidad y está prohibida por el derecho internacional.
        \end{itemize}

    Existen disitnos posibles escenarios para los casos de personas desaparecidas según Charles Twardy: avalancha, criminal, abatido, evadido, investigación, perdido, médico, ahogamiento, retrasado, varado o atrapado y trauma.  

    


    \subsection{Dispositivos de búsqueda}
    Un dispositvio de búsqueda es multidisciplinar, como la ingeniería biomédica, aunque normalmente los casos de desapariciones los llevan a cabo FFCC de seguridad o bomberos coordinados por el 112, pueden intervenir distintso grupos desde bomberos, guardia civil, policía protección civil, sanitarios, agentes de medioambiente, organizaciones de voluntarios... \cite{ManualDeBusquedayRescate}

    Todos los equipos durante una búsqueda se encuentran comunicados con el MANDO.

    Se suelen realizar difusiónd e alertas por redes sociales y/o medios de comunciación. También el sistema es-alert. 

    Normalmente todos los participantes en una búsqueda se encuentran divididos en equipos de búsqueda y en cada equipo habrá un jefe de equipo que cuenta con formación específica en búsqueda de personas desaparecidas.

    En una búsqueda podemos encontrar 3 perfiles diferentes:
    \begin{itemize}
        \item Buscador
        \item Jefe de equipo
        \item Planificador
    \end{itemize}

    

%%%%%%%%%%%%%%%%%%%%%%%%% CONOCIMIENTOS PREVIOS %%%%%%%%%%%%%%%%%%%%%%%%%
    \section{Conocimientos previos II. GPS Cerebral.}\label{previos}
    \textcolor{red}{LA MAYOR PARTE DE LA INFO DE AQUÍ ES DEL PRIMER MANUAL DE BÚSQUEDA Y RESCATE DE PERSONAS DESAPARECIDAS EN CASTELLANO.} \cite{ManualDeBusquedayRescate}

    \textcolor{red}{ Sentar las bases teóricas y técnicas en las que se apoya la aplicación.}

    Antes de adentrarnos en el el estado del arte y el desarrollo de la aplicación, nos adrentarmos en los fundamentos biológicos y teóricos en los que nos apoyaremos para el desarrollo de nuestro trabajo.

    \subsection{GPS cerebral}
    Se puede decir que la orientación es el resulatado de la actividad neuronal de regiones específicas.

    \textcolor{red}{Buscar un poco más del premio Nobel para confirmar que la frase está correctamente dicha y ampliar información} \cite{GPSCerebral_MOSER,GPS_Cerebral_O_KEEFE}

    El GPS Interno del que constan todas las personas fue descubierto y definido por los premios Nobel John O'Keele, May-Britt Moser y Eduard I.Moser. El GPS interno es el sistema de posicionamiento que tiene el cerebro para conocer las posiciones de la persona y de otros objetos o lugares. En el GPS interno se ven implicados los sentidos y dos zonas del cerebro, el hipocampo donde formamos los mapas y la corteza entorrinal donde se tomarán las decisiones, que se encuentra muy cercana al hipocampo. 

    \textcolor{red}{El mapa cognitivo se trata de una representación del entorno que se genera en el hipocampo, se codificarían señales visuales, auditivas y somato sensoriales.}
    
    \textcolor{red}{Incluir una imagen del cerebro donde se encuentre marcado el hipocampo y la corteza entorrinal}

    El hipocampo forma parte del sistema límbico (hipocampo+amígdala), que se trata del sustrato morfológico de la memoria, aprendizaje, motivación y las emociones. 

    En el mapa cognitivo también influye la memoria declarativa, es decir, nuestros recuerdos. 
    
    \textcolor{red}{(La memoria declarativa se divide en episodica y semántica.)} 

    Hay estudios que demuestran que si una persona se encuentra en constante uso de la orientación, como los taxistas, tienen un hipocampo más activo que las personas que no hacen tanto uso de la orientación. También existe una correlación entre la buena orientación y la cantidad de sustancia gris en la zona.

    El hipocampo se divide en 4 regiones CA1, CA2, CA3 y CA4. \cite{AnatomiaHipocampo}
    %https://www.kenhub.com/es/library/anatomia-es/hipocampo

    En estos tejidos encontramos 4 tipos de células implicadas en el GPS Interno. Las células de lugar, las células de red, las células de dirección y las células de límite.
    
    \subsubsection{Células nerviosas de lugar o de posicionamiento}
    Se trata de neuronas que se encuentran en el hipocampo que responden selectivamente al lugar donde se encuentra la persona. Forman sistemas de navegación alocéntricos de detección entre objetos. Por decirlo de otra forma, se podría decir que se trata de los puntos de refernecia que utilizamos para saber donde nos encontramos.
    

    Estas células se activan cuando se está en un lugar específico.

    Dan información sobre la posición de objetos alrededor, de partes del propio cuerpo o de lugares y permiten un desplazamiento sin errores.

    Se cree que las células de lugar tambien son las células de tiempo, aquellas que codifican la percepción del tiempo. Se desconoce su potencial de acción.

    \subsubsection{Células de red}
    Son neuronas que se encuentran en el córtex entoriinal. Se llaman células de red porque se organizan en patrones de células en hexágonos o triángulos. Actualmente su uso es desconocido pero se sabe que se encuentran relacionadas con las células de lugar. Se cree que podrían funcionar cómo la triangulación de telefonía, para poder conocer las coordenadas. 

    La combinación del mapa cognitivo, las células de lugar y estas mismas permitirían crear el mapa cerebral.

    \subsubsection{Células de dirección}
    Estas neuronas permitirían conocer la dirección en la que se mueve el sujeto. 
    

    \subsubsection{Células de límite}
    Son neuronas que se encuentran en el subículum del hipocampo. Estas neuronas experimentan una mayora activación cuando el sujeto se encuentra cerca de algún límite, cómo puede ser una pared o un precipicio.

    
    
    
    Existen numerosas enfermedades y trastornos que pueden afectar a nuestro GPS cerebral como la hipoglucemia, hiperglucemia, hipotermia, hipertermia, demencia, enfermedad mental, discapacidad mental y autismo.

    Por ejemplo, con la hipoxia se disminuye la irrigación a esas zonas del cerebro. Y con la edad se ve afectada la memoria a corto plazo por lo que para intentar mantenerla hay que ejercitarla lo máximo posible.
    

    \subsection{GPS cerebral: Alzheimer}
    Se ve afectado el hipocampo y la corteza entorrinal, de hecho es una de las primeras estructuras del cerebro que se deteriora, por lo que se produce una pérdida del sentido de la orientación, como uno de los primeros síntomas. Su causa es la acumulación de proteínas tau y de placas de proteína beta-amiloide que destruyen las neuronas. 
    
    En concreto, unas de las primeras zonas que se ven afectadas son la corteza entorrinal por la acumulación de proteína Tau y la región CA1 por la acumulación de péptido amiloide que provoca poco a poco la desconexión entre las neuronas, por lo que se va perdiendo memoria a corto plazo y existe una deficiencia en la activación de las células de red. 
    
    Es una de las enfemedades con más incidencia y por el momento no tiene cura.
    
    La enfermedad de Alhzéimer se caracteriza por una memoria deteriorada, dificultad para calcular el tiempo y las distancias. En definitiva, de una pérdida de las habilidades viso-espaciales. Actualmente no existe cura alguna.

    Las personas con Alzhéimer avanzado son incapaces de crear un mapa cognitivo, se encuentran perdidas todo el tiempo y su orientación y puntos de referencia se basan en todo lo que vivieron en el pasado, en la información que tienen almacenada a largo plazo.
    
    

    \subsection{GPS cerebral: Otras}
    Personas con demencia, discapacidad intelectual o enfermedad mental tienen disminuida también la capacidad de generar mapas cofnitivos.

    Sin embargo cabe destacar que las personas con discapacidad intelectual tienen una mejor capacidad que las personas cuya edad se les ha atribuido puesto que cuentan con la años de experiencia. Principalmente se basan en rutinas que si llegan a cambiar pueden impedir que consturyan impediatamente el nuevo mapa cognitivo.

    Los niños todavía no son capaces de generar mapas cognitivos completos puesto que no han desarollado completamente esa habilidad. Normalmente se empieza a desarrollar a partir de los 8-9 años. De hecho los niños pequeños tampoco suelen tener metaconocimiento de sus habilidades espaciales, por lo que pueden perderse con mayor facilidad.

    \subsection{GPS cerebral: Emociones}
    Como se ha mencionado antes, las emociones también juegan un papel principal en la orientación de las personas.

    Por ejemplo, una persona con una alta excitación percibe menor cantidad de señales ambientales, mientras que una persona con poca excitación puede incurrir en pensamientos difusos. Por lo que como en todo en esta vida, se busca el equilibrio.

    Y no solo la excitación, el miedo también influye, interfiere en el funcionamiento normal mental, en la concentración y en la capacidad de resolver problemas. Además, una persona que siente miedo busca la auto-preservación por lo que se produce una mayor adrenalina y mayor cantidad de sangre en las piernas por lo que su impulso normalmente será moverse.

    Por otro lado, también se debe tener en cuenta si una persona desaparece sola o en grupo, puesto que las personas actúan diferente y suelen sentir menos miedo y siguen comportamientos más racionales. De hecho, lo normal es que si desaparece un grupo estas personas se queden en grupo o algunas vayan a buscar ayuda.



    
%%%%%%%%%%%%%%%%%%%%%%%%%%%%% ESTADO DEL ARTE %%%%%%%%%%%%%%%%%%%%%%%%%%%%%
    \section{Estado del arte}\label{estadoArte}
    %\textcolor{red}{Revisión exhaustiva y crítica de la literatura existente y las tecnologías actuales en un campo específico. En el contexto de una app, el estado del arte implica analizar las innovaciones, tendencias, metodologías, y las soluciones más avanzadas que existen hasta la fecha en ese ámbito particular. Su objetivo es comprender lo que ya se ha logrado y los desafíos aún por resolver, para identificar áreas de mejora e innovación.}
    
    En la actualidad existen numerosas herramientas y sistemas especializados que se utilizan en operaciones de búsqueda y rescate de personas desaparecidas. En este apartado se presenta una revisión de soluciones relevantes, incluyendo algunas que, si bien no han sido concebidas específicamente para este ámbito, podrían aportar funcionalidades valiosas si se adaptan adecuadamente. A partir de esta recopilación, se realizará un análisis comparativo para poder identificar las características más interesantes para el diseño de la aplicación propuesta en este trabajo. 
    
    \subsection{Herramientas y sistemas específicos de casos de personas desaparecidas}
    \textcolor{red}{Incluir aquí las herramientas, alertcops, desaparecido,... Todo menos Alerthor y alguna de las otras aplicaciones tipo mapas que son de agricultura etc... QUIZÁS ORDENARLAS EN FUNCIÓN DE SU USO EN LA BÚSQUEDA, PRIMERO LA MONITORIZACIÓN DE LAS PERSONAS Y LAS DENUNCIAS CON ALERCOPS, LUEGO SISTEMAS DE BASES DE DATOS, LOS SISTEMAS DE ALERTAS Y LOS SISTEMAS DE BÚSQUEDA/MAPAS.} 
    
    En esta sección se presentan soluciones tecnológicas desarrolladas específicamente para la localización de personas desaparecidas. Se incluyen aplicaciones, plataformas, dispositivos y metodologías, con énfasis en sus funcionalidades, limitaciones y grado de adopción por parte de organismos especializados.

    \subsubsection{Alertcops}
    AlertCops representa una de las herramientas más relevantes en el ámbito de la seguridad ciudadana en nuestro país, pese a no ser suficientemente conocida por la población. Se trata de una aplicación móvil gratuita desarrollada por el Ministerio del Interior español, cuyo objetivo principal es reforzar la seguridad ciudadana mediante una comunicación directa, discreta y rápida con los cuerpos de seguridad del Estado, como la Guardia Civil y la Policía Nacional. \cite{Alertcops} 

    La aplicación permite a la ciudadanía informar sobre delitos, incidencias o situaciones sospechosas, así como compartir datos con las autoridades. En el contexto de personas desaparecidas, una de las funcionalidades más importantes es el modo guardián. Esta funcionalidad permite compartir la ubicación en tiempo real con familiares o directamente con los cuerpos de seguridad, lo que facilita una localización rápida y precisa en caso de emergencia.

    Además, Alertcops cuenta con un chat en tiempo real con la Policía y la Guardia Civil, lo que permite recibir orientación inmediata en casos de desorientación o riesgo. También incluye un botón de socorro (SOS) que, al ser activado, envía automáticamente una alerta a las autoridades acompañada de un breve audio y la geolocalización del usuario. \cite{alertcops_folleto}
    

    \subsubsection{Aplicación Desaparecido}

    \textcolor{red}{contactar con el CEO (Eduardo) para preguntar más info o que ha pasado porque no hay ninguna noticia después de febrero de 2023 (bueno en linkedin tienen una única publicación de hace 10 meses), y creo que es algo muy muy parecido a mi idea}
    
    Desaparecido es un software SAR (Search and Rescue) orientado al apoyo de los equipos de emergencia en la búsqueda de personas desaparecidas\cite{webDesaparecido, Github_DesaparecidoApp}, con el propósito de mejorar la eficiencia y coordinación de las operaciones. Entre sus funcionalidades destaca el monitoreo y control en tiempo real de los datos asociados a cada caso, la documentación de los planes de actuación, y la coordinación entre los distintos equipos implicados. Esta herramienta se basa en procesos estandarizados que facilitan el trabajo conjunto entre cuerpos de emergencia, y garantiza una gestión segura de la información, priorizando la privacidad de los datos personales. 

    Actualmente, este software se encuentra en fase de desarrollo y solo está disponible una demo. Su propuesta consiste en digitalizar los procedimientos que, en la mayoría de los casos, se realizan de forma manual, ofreciendo una plataforma integral para la gestión de casos de personas desaparecidas. 
    
    Según la información disponible en redes sociales y en su sitio web, se están incorporando funcionalidades avanzadas basadas en Machine Learning como el cálculo de la vulnerabilidad del caso, que permite estimar el grado de urgencia mediante parámetros como la edad de la persona desaparecida, su estado físico o las características del terreno. Asimismo, se busca desarrolla un sistema de predicción de zonas de búsqueda, basado en modelos que identifican las áreas con mayor probabilidad de localización de la persona desaparecida. Por otro lado, desde la aplicación también se contempla la gestión de voluntarios, facilitando la organización de los recursos humanos disponibles, así como la comunicación en tiempo real con los familiares, ofreciendo actualizaciones constantes sobre el estado de la búsqueda.
    

    \subsubsection{Win CASIE III}

    Win CASIE III es un software SAR (Search and Rescue) gratuito\cite{WCASIE3}, especialmente utilizado en operaciones de búsqueda de personas desaparecidas en entornos rurales o naturales y recomendado desde el Manual de Búsqueda y Salvamento Terrestre\cite{ManualDeBusquedayRescate}. Su diseño está orientado a apoyar a los responsables del operativo durante todas las fases de la intervención, desde la planificación inicial hasta la elaboración del informe final, proporcionando distintas herramientas para la toma de decisiones y la documentación del proceso.

    En la fase inicial de la búsqueda, el software permite acceder a una amplia variedad de formularios, como los relativos a la identificación de la persona desaparecida, la descripción del incidente, la evaluación de la urgencia del caso o la recogida sistemática de pistas. También ofrece la posibilidad de utilizar formularios oficiales de instituciones como el Emergency Management Institute (EMI), así como información sintetizada sobre el comportamiento de personas desaparecidas (Lost Person Behavior), que orienta las líneas de investigación en función de las circunstancias del caso. Además, Win CASIE III facilita el registro detallado de las decisiones tomadas y toda la información, inlcuyendo datos sobre la persona desaparecida, el nivel de urgencia, los recursos disponibles y requeridos, la estructura del operativo y la composición de los equipos implicados.

    Durante la fase de búsqueda sobre el terreno, el software permite imprimir formularios de consenso para expertos, calcular el consenso general a partir de sus respuestas y almacenarlo como parte del informe. También se pueden generar hojas de asignación de equipos, informes intermedios y otros documentos operativos. Asimismo, Win CASIE III incorpora herramientas para calcular parámetros clave como la probabilidad de detección (POD) o la probabilidad de área (POA), además de permitir la segmentación del terreno en zonas según las necesidades del operativo. La plataforma facilita el registro de pistas encontradas, la realización de análisis hipotéticos tipo What if?, y ofrece recomendaciones sobre la distribución óptima de los recursos para maximizar la probabilidad de éxito. Todo el proceso queda almacenado, lo que permite mantener una trazabilidad completa de las decisiones y acciones ejecutadas.

    Por último, una vez finalizado el operativo, el software permite generar un informe final completo que recoge de forma estructurada todas las acciones realizadas, las decisiones tomadas y los resultados obtenidos.
    


    \subsubsection{Sistema ES-ALERT}

    ES-Alert es el sistema español de alertas públicas, desarrollado en el cumplimiento del artículo 110 de la Directiva Europea 2018/1972\cite{DirectivaUE_EsAlert}, que exige a los Estados miembros la implementación de mecanismos capaces de proporcionar instrucciones claras y oportunas a la población en situaciones de emergencia mediante telefonía móvil. En España, este sistema se integra en la Red de Alerta Nacional de Protección Civil como complemento a los canales tradicionales de alerta y está gestionado por el Ministerio del Interior\cite{SistemaEsAlert,ESAlertProteccionCivil}.
    
    Su funcionamiento se basa en la tecnología Cell Broadcast, que permite enviar mensajes de forma simultánea e indiscriminada a todos los teléfonos móviles conectados a las antenas de una zona geográfica determinada. Por tanto, su eficacia depende de que los dispositivos estén encendidos y con cobertura. Para emitir una alerta, se configuran distintos parámetros como el área de afectación de la emergencia, la hora de inicio y duración de la alerta, el nivel de criticidad y el contenido del mensaje. Aunque habitualmente se trata de mensajes de texto, existe la posibilidad de incluir imágenes. 

    El sistema ES-Alert está disponible para los centros de coordinación de emergencias de las comunidades autónomas y para el Centro Nacional de Seguimiento y Coordinación de Emergencias. Este sistema se ha empleado en situaciones como durante fenómenos meteorológicos extremos o movimientos sísmicos, donde la rapidez en la difusión de la información es esencial para proteger a la población. Por otro lado, se ha de mencionar que cumple con las normativas de privacidad nacionales y europeas, ya que no almacena datos personales de los receptores de las alertas.
    
    Aunque no ha sido diseñado específicamente para casos de desaparición, sí es una herramienta de apoyo en operaciones de búsqueda y rescate, siempre y cuando no se comprometa la seguridad de la persona desaparecida. En estos casos, ES-Alert sirve para difundir información clave a la población localizada en el área de búsqueda, facilitando la colaboración ciudadana y ampliando el alcance de los equipos de emergencia\cite{ESAlertCastillaLeon2025}. Por ello, las características de este sistema pueden ser muy interesantes para su adopción en nuestra aplicación o incluso que en algún momento se encuentren interconectados.

    \subsubsection{Bases de datos de personas desaparecidas: Sistema PDyRH y Database ISRID}

    \textcolor{red}{Dar Vuelta/Revisar}
    
    \textcolor{red}{Citar el plan estrategico en materia de personas desaparecidas.} \cite{plan_desaparecidos_2022}
    
    Sistema PDyRH: sistema de Personas Desaparecidas y Restos Humanos sin Identificar (PDyRH). Es el empleado por las Fuerzas y Cuerpos de Seguridad en la investigación de desapariciones. Es coordinado por el CNDES (Centro Nacional de Desaparecidos) que además colabora con las Administraciones Públicas y otras instituciones y organizaciones públicas y privadas habilitadas, tanto nacionales como internacionales en la difusión de avisos, alertas o peticiones de colaboración a la población. https://www.interior.gob.es/opencms/en/detail-pages/article/Las-Fuerzas-y-Cuerpos-de-Seguridad-registraron-26.003-denuncias-por-desaparicion-de-personas-en-2022/



    \subsubsection{Sistemas basados en Inteligencia Artificial: DroneFinder y RES-DesVi}

    \textcolor{red}{Dar Vuelta/Revisar}
    
    DroneFinder: https://itg.es/dronefinder-plenamente-operativaermite-localizar-personas-desaparecidas-con-drones-ia/

    Respecto a los sistemas de predicción en este ámbito, son escasos pero en un tiempo donde la inteligencia artificial está a la orden del día el desarrollo de sistemas de predicción irá creciendo en todos los ámbitos.... 

    De hecho ya hay investigaciones sobre sistemas de predicción para el apoyo en las búsquedas de personas desaparecidas. Una des las herramientas que se está desarrollando es SER-DesVi que se trataría de un sistema predictivo de valoración del riesgo de lesión o fallecimiento por causas suicidas u homicidas. \cite{SERDesVi, SERDesVi_Validacion}

    SER-DesVi se basa en distintos factores de riesgo y protección para realizar las predicciones. También tiene como objetivo es que sirva como herramienta estandarizada de recolección de datos sociodemográficos, psicosociales y criminológicos de desapariciones en España. 

    Este estilo de herramientas de predicción favorecen a la priorización de líneas de investigación ante una desaparición y a la mejora de gestión de los recursos de los operativos.

    Sin embargo, falta una inclusión en los sistemas multirespuesta para que sea realmente efectiva la toma de decisiones en los operativos de búsqueda de personas desaparecidas basada en la evidencia. 
    
    \subsection{Otras herramientas y sistemas}
    Además de las soluciones diseñadas para este fin, existen herramientas provenientes de otros ámbitos (como la gestión de emergencias) que podrían ofrecer ventajas significativas si se adaptan al contexto de búsqueda de personas desaparecidas. Esta subsección explora dichas tecnologías con el objetivo de ampliar el espectro de posibilidades para el diseño de nuestra aplicación.

    \subsubsection{Mapas: SIGPAC y IGME}

    \textcolor{red}{Dar Vuelta/Revisar}

    SIGPAC: mapas https://sigpac.mapa.es/fega/visor/
    
    IGME: mapas https://info.igme.es/cartografiadigital/portada/default.aspx?mensaje=true

    \subsubsection{Alerthor}

    \textcolor{red}{Dar Vuelta/Revisar}
    
    Alerthor es una plataforma de gestión de alertas emergencias y monitoreo geográfico diseñado para equipos de rescate, bomberos y otros servicios de seguridad. Desde esta plataforma se pueden coordinar recursos en tiempo real y geolocalizar tanto a personas como a vehículos desde un único monitor. Su sistema de alertas funciona incluso sin conexión a internet. Además, permite personalizar mapas y generar informes de los operativos. \cite{webAlerthor, appAlerthor}
        
    Aunque no se trate de una herramienta especializada en operativos de personas desaparecidas, es una herramienta que podría ser muy útil para organizar y coordinar las búsquedas. Por ello, muchas de las funcionalidades disponibles las tendré en cuenta para definir nuestra aplicación especializada en la búsqueda y rescate de personas desaparecidas. 

    


    \subsection{Análisis comparativo}
    A partir de la revisión anterior, se realiza un análisis comparativo que permite identificar patrones, fortalezas y debilidades entre las distintas herramientas. Este ejercicio busca extraer los elementos más relevantes y transferibles que puedan servir como base para la definición funcional y técnica de la aplicación que se pretende desarrollar.
    
    \textcolor{red}{Análisis Comparativo:}

    \textcolor{red}{Realizar una comparación entre los distintos proyectos (características, ventajas y problemas de las soluciones existentes.)}
    
    Los principales problemas y ventajas que ofrecen las soluciones actuales son:
    (Crear una tabla)

    Existen sistemas de alerta ciudadana y redes sociales especializadas como:
    \begin{itemize}
        \item es-alert
        \item Web del CNDES
        \item asociaciones como sosdesaparecidos
    \end{itemize}
    
    \textcolor{red}{Análisis de métodos y tecnologías:}
    
    %\textcolor{red}{- Revisa tecnologías disruptivas que se estén utilizando en la coordinación de emergencias (uso de drones, algoritmos de inteligencia artificial para análisis de datos, etc.).}
    
    \subsection{Sistemas de predicción actuales}

    \textcolor{red}{INCLUIDO PREVIAMENTE EN EL APARTADO DE SISTEMAS BASADOS EN INTELIGENCIA ARTIFICIAL}

    Respecto a los sistemas de predicción en este ámbito, son escasos pero en un tiempo donde la inteligencia artificial está a la orden del día el desarrollo de sistemas de predicción irá creciendo en todos los ámbitos.... 

    De hecho ya hay investigaciones sobre sistemas de predicción para el apoyo en las búsquedas de personas desaparecidas. Una des las herramientas que se está desarrollando es SER-DesVi que se trataría de un sistema predictivo de valoración del riesgo de lesión o fallecimiento por causas suicidas u homicidas. \cite{SERDesVi, SERDesVi_Validacion}

    SER-DesVi se basa en distintos factores de riesgo y protección para realizar las predicciones. También tiene como objetivo es que sirva como herramienta estandarizada de recolección de datos sociodemográficos, psicosociales y criminológicos de desapariciones en España. 

    Este estilo de herramientas de predicción favorecen a la priorización de líneas de investigación ante una desaparición y a la mejora de gestión de los recursos de los operativos.

    Sin embargo, falta una inclusión en los sistemas multirespuesta para que sea realmente efectiva la toma de decisiones en los operativos de búsqueda de personas desaparecidas basada en la evidencia. 

    


    
    
    \subsection{Resumen de identificación de huecos}
    \textcolor{red}{Identificación de huecos en la literatura:}
    
    \textcolor{red}{- Señala las áreas no suficientemente cubiertas o las limitaciones de las actuales herramientas, lo que justifica tu propuesta de solución.}
    
    \textcolor{red}{Huecos identificados:}
    Sin embargo, la comunicación entre los grupos de rescate, profesionales en caso de necesitar la colaboración ciudadana es bastante limitada, utilizando redes sociales y contactos conocidos. Por lo que la creación de una app que aumente la esa colaboración y comunicacion entre los profesionales y los civiles dispuestos a aportar su granito de arena, con el fin de disminuir tiempos de búsqueda y aumentar las posibilidades de encontrar con vida a la persona desaparecida.

    Además, en un mundo donde cada vez más tecnológico y donde tenemos a la orden del día la IA (inteligencia artificial), no se están aprovechando todos los recursos. Por ello también se intenta incluir en nuestra solución una funcionalidad de predicción para poder enfocarse en las zonas donde sería más probable encontrar a la persona desaparecida en función de sus características. (Edad, enfermedades...)
    
    %\textcolor{red}{- Señala cuáles son los principales desafíos del procedimiento tradicional (por ejemplo, la dispersión de información, la lentitud en la coordinación, problemas de actualización en tiempo real) y cómo tu app puede suponer una mejora.}
    
    
    

    
%%%%%%%%%%%%%%%%%%%%%%%%%%%%%%%% OBJETIVOS %%%%%%%%%%%%%%%%%%%%%%%%%%%%%%%%
    \section{Objetivos}\label{obj}
    %\textcolor{red}{Objetivos generales y específicos: Define qué se pretende lograr con el desarrollo de la app, tanto a nivel funcional (gestión, búsqueda y coordinación) como en términos de impacto social.}
    
    \textcolor{red}{En este trabajo, se parte de la hipótesis de que \emph{.....}.}
    
    De acuerdo con ello, el objetivo principal de este trabajo es el de conseguir una \textbf{aplicación de apoyo en la búsqueda y rescate de personas desaparecidas}. Para conseguir dicho fin, son necesarios los subobjetivos que se detallan a continuación:
    
    \begin{itemize}
        \item \textbf{Subobjetivo 1:} Investigar los procedimientos y las herramientas actuales de búsqueda de personas desaparecidas.
        
        \item \textbf{Subobjetivo 2:} Definir las características de la aplicación.
        
        
        \item \textbf{Subobjetivo 3:} Definir las gestiones internas de la aplicación.
        
        
        \item \textbf{Subobjetivo 4:} Implementar un prototipo de la aplicación.

        \item \textbf{Subobjetivo 5:} Sentar las bases de un sistema básico de predicción de la localización de una persona desaparecida en base a enfermedades u otras características de la persona desaparecida.
        

    \end{itemize}
    
%%%%%%%%%%%%%%%%%%%%%%%%%%%%%%% DESCRIPCIÓN %%%%%%%%%%%%%%%%%%%%%%%%%%%%%%%
    \section{Descripcion general de la solucion}\label{descripcion1} % COMENTAR SI NO INTERESA
    Se pretende realizar una app de apoyo en la búsqueda de personas desaparecidas. Esta app se basará en un sistema de alertas, predicción y geolocalización. Pretende reunir a la mayor cantidad de voluntarios posible y llevar un control de los asistentes a la búsqueda. Además, servirá también para logística de los equipos profesionales de búsqueda y optimización de los recursos con el fin de encontrar a la persona desaparecida en el menor tiempo posible.

    La app estaría disponible no solo para grupos profesionales o voluntarios profesionales, sino que también para cualquier persona interesada en ayudar. Estos civiles podrían gestionar el área donde se quiere recibir las alertas y aceptar o denegar ir a ayudar en caso de alerta de búsqueda en la zona. Durante la búsqueda se mantendrá informados y geolocalizados a los asistentes.

    Respecto a la ayuda logística de la búsqueda, se podrán definir los distintos perfiles necesarios o las tareas a realizar, colocar puntos de control y ofrecer retroalimentación basada en estadísticas de desapariciones anteriores y la información del desaparecido, como puede ser edad o enfermedades, por ejemplo, si desaparece en la montaña una persona de mediana edad con diabetes la apliación indicaría que algunos de los puntos de control deben de ser zonas con alguna fuente de agua ya que está demostrado que las personas con diabetes suelen tener mucha sed y por supervivencia se acercarán a esas fuentes de agua.
    
    Se ha utilizado la herramienta Github \cite{GitHub} para llevar un control de todo el proyecto. Se puede acceder a la información del mismo en el siguiente repositorio: ...

    \textcolor{red}{REVISAR }

    La app de búsqueda y rescate de personas desaparecidas que se presenta en este trabajo requiere de la gestión y análisis de datos de personas desaparecidas, grupos de rescate, voluntarios e instituciones, y en el caso de personas desparecidas también se puede trabajar con sus datos médicos como información de enfermedades o trastornos. Como la mayoría de datos que se van a tratar en la app son datos sensibles, como información médica y personal, requiere tener muy en cuenta la protección de esos datos y la privacidad. \textcolor{red}{TRATAMIENTO DE LOS DATOS}

    La automatización de alertas en aplicaciones de seguridad y de colaboración ciudadana, pueden transformar la forma en que los organismos de rescate gestionan los recursos y aumentan las probabilidades de éxito en situaciones críticas para poder salvar vidas. \textcolor{red}{FUERA/ EXPLICACIÓN DE ALERTAS}

    El desarrollo de la app pretende mejorar la gestión y coordinación de  los operativos de búsqueda de personas desaparecidas. 
    Proporcionando a los profesionales herramientas que faciliten su trabajo. Además, gracias a la herramienta que se pretende desarrollar se promueve la colaboración entre los distintos organismos, grupos de rescate, asociaciones y civiles en general. \textcolor{red}{DESCRIPCIÓN DE LA SOLUCIÓN}
    
    Alcance y limitaciones: Describe brevemente qué aspectos cubrirá el proyecto y establece los límites de la misma investigación y desarrollo.\textcolor{red}{DESCRIPCIÓN DE LA SOLUCIÓN}

    Probabilidad del área, la probabilidad de encontrar a la persona desaparecida en ese área. Existen diferentes métodos para calcularla

\begin{comment}
    \textcolor{red}{Este es el núcleo del trabajo, donde describes en detalle tu propuesta:}
    
    \textcolor{red}{Arquitectura del sistema:}

    
    \textcolor{red}{Funcionamiento y funcionalidades:}
    
    \textcolor{red}{- Gestión de casos: Registro de incidentes y seguimiento del estado del caso.}
    
    \textcolor{red}{- Búsqueda y localización: Integración con mapas, geolocalización en tiempo real, uso de algoritmos heurísticos para determinar zonas críticas.}
    
    \textcolor{red}{- Coordinación: Herramientas de comunicación en tiempo real (chat, notificaciones push, llamadas de emergencia) para coordinar a los equipos de búsqueda y coordinación.}
    
    \textcolor{red}{- Interfaz de usuario (UI/UX): Describir los principales elementos de la interfaz y la experiencia de usuario, incluyendo la accesibilidad y la claridad en la visualización de datos críticos.}
    
    \textcolor{red}{Tecnologías y frameworks empleados:}
    
    \textcolor{red}{- Lista y justifica la elección de lenguajes, frameworks y bases de datos.
    
    \textcolor{red}{- Menciona si utilizas metodologías de desarrollo ágil y cómo integras las pruebas (QA, testing de usabilidad, etc.)}

    \textcolor{red}{Comparativa con Procedimientos Actuales:}

    \textcolor{red}{- Al describir la arquitectura y funcionalidades de la app, puedes incluir una sección comparativa que muestre cómo los pasos de tu solución se alinean o mejoran cada fase de los procesos tradicionales.}
    
    \textcolor{red}{- Utiliza diagramas de flujo o infografías para visualizar el procedimiento actual versus el procedimiento propuesto mediante la app. Esto no solo aporta claridad, sino que también destaca la innovación que representa tu solución.}

    

    \subsubsection{Tipos de información a tratar}
    \textcolor{red}{Tipo de información con la que se trata en las desapariciones.}
    % Guía breve para familiares y organizacioens de la sociedad civil sobre la recolección de información sobre personas desaparecidas. Comisión Internacional sobre Personas Desaparecidas (ICMP)
    
    En los casos de personas desaparecidas se trata con distintos tipos de información en concreto podemos dividir entre\cite{GuiaFamiliaresICMP}:
    \begin{itemize}
        \item Información personal de la persona desaparecida: incluye datos básicos como nombre, fecha de nacimiento, características físicas y última localización conocida.
        \item Información acerca de la desaparición: fecha, lugar o circunstancias asociadas a la desaparición.
        \item Información personal de la persona que reporta la desaparición: nombre e información de contacto.
        \item Información personal de familiares: nombres, información de contacto y relación de parentesco.
        \item Información de personal de búsqueda o voluntarios.
    \end{itemize}

     

    \textcolor{red}{Bases tecnológicas:}
    
    \textcolor{red}{- Introduce las tecnologías móviles y frameworks de desarrollo (Android, iOS, híbridos, etc.).}
    
    \textcolor{red}{- Habla sobre sistemas de localización y mapeo (GPS, GIS) y su fiabilidad.}
    
    \textcolor{red}{- Si aplica, menciona aspectos relacionados con IoT o sensores biomédicos que puedan integrarse para aspectos relativos a la salud o localización.}
    
    \textcolor{red}{Fundamentos en ingeniería biomédica:}
    
    \textcolor{red}{- Expón cómo los principios biomédicos (por ejemplo, monitorización de signos vitales en situaciones de estrés o trauma) podrían aportar valor al sistema.}
    
    \textcolor{red}{- Referencias a estudios previos o metodologías biomédicas utilizadas en emergencias y rescates.}
    
    \textcolor{red}{Normativa y aspectos éticos:}
    
    \textcolor{red}{- Breve revisión de normativas sobre protección de datos (GDPR u otras) en el manejo de información sensible.}
    
    \textcolor{red}{- Consideraciones éticas en el uso y difusión de información de personas desaparecidas.}

    \textcolor{red}{Protocolos de Actuación en Emergencias:}

    \textcolor{red}{- Profundiza en los métodos y procedimientos que siguen las autoridades y organismos oficiales.}
    
    \textcolor{red}{- Incluye referencias a normativas, estudios de caso y protocolos oficiales (por ejemplo, documentos de políticas de emergencias o guías de actuación de instituciones).}
    
    \textcolor{red}{Este enfoque ayuda a asentar las bases teóricas y técnicas sobre las cuales se apoya el desarrollo de la aplicación y resalta la relevancia de integrar, adaptar o mejorar lo que ya se aplica en la práctica.}

\end{comment}
    
    Durante el trabajo he tenido que aprender JavaScript y Node.js para poder implementar el código de la aplicación híbrida en React Native.

    He consultado detalles de las caracteristicas de la aplicación y la usabiliad de la aplicación con el grupo GREM (Grupo de Rescate Espeleológico y de Montaña) de Burgos.
    
